% \resultpct{<code>_better}: the "k of n" macro as a rounded percentage.
\ExplSyntaxOn
\NewDocumentCommand{\resultpct}{m}{
  \cs_if_exist:cTF {opera@res@#1}
    {
      \tl_set:Nx \l_tmpa_tl { \use:c {opera@res@#1} }
      \tl_remove_all:Nn \l_tmpa_tl {,}
      \seq_set_split:NnV \l_tmpa_seq {of} \l_tmpa_tl
      \resultstyle{ \fp_eval:n { round(100 * (\seq_item:Nn \l_tmpa_seq {1})
        / (\seq_item:Nn \l_tmpa_seq {2}), 0) } \% }
    }
    { \placeholder{X\%} }
}
\ExplSyntaxOff

% OPERA manuscript body, consolidated edition (30 October 2026).
% Voice and structure of opera_body.tex (rehearsal review of 29 October), held
% to the Lancet Digital Health rules: Summary <= 300 words in five headed
% paragraphs, Research in context without citations, main text <= 3,500 words,
% <= 30 references. The thread through every section is borrowing: whether to
% borrow, what to borrow and how to borrow. Flat and disease-balanced summaries
% are co-reported (lock amendment, 30 October 2026).
\begin{abstract}
\textbf{Background} Prognosis shapes treatment, supportive care and patients'
decisions about their future, yet most prognostic tools learn from one disease
and one outcome at a time. We asked whether, what and how to borrow information
across diseases, outcomes and patient histories to predict risk in lymphoid
cancers, a specialty of many uncommon diseases.

\textbf{Methods} We studied 37,531 Danish adults starting first-line treatment
for 24 lymphoid cancers, predicting 69 outcomes at six horizons from 30 days to
two years from routinely recorded histories. Otherwise identical sequence models
learned from one disease--outcome pair, one outcome across all diseases, all
outcomes within one disease, or every disease and outcome together. \opera{}, a
single model for all tasks, adds corrections specific to each patient's disease
and outcome and lets each prediction borrow those of other tasks. Encoders were
pretrained on the cohort or the national population. Comparators were per-task
and pooled XGBoost and the prognostic indices in clinical use, in patients
treated from 2023.

\textbf{Findings} Borrowing from other diseases increased the area under the
ROC curve (AUC) by \result{P1_auc}{X} over single-task models. Adding all other
outcomes to the disease-pooled model changed AUC by \result{P4_auc}{Y}
(\placeholder{claim label}), so one model could predict every outcome without
loss. \opera{} improved on the joint model by \result{R1_auc}{X}, with borrowed
corrections adding \result{router_borrow_AB_auc}{Z}, and differed from per-task
XGBoost by \result{O1_auc}{Y}. \placeholder{National pretraining.}
\placeholder{Against the N prognostic indices in routine use for these diseases,
for each index's own endpoint and patients: the number in which OPERA was
better, with the NCCN-IPI result in DLBCL.}

\textbf{Interpretation} \placeholder{From the test results: borrowing across
diseases improved prognosis most where data were scarce, and one model served
every disease and outcome at no cost. How this compares with tuned tabular
models and clinical indices, without claiming clinical benefit.}

\textbf{Funding} \placeholder{Funders, to be confirmed by the authors.}
\end{abstract}

\noindent\fbox{\parbox{\dimexpr\textwidth-2\fboxsep-2\fboxrule}{%
\small\textbf{Research in context}\par
% No citations in this panel.
\textbf{Evidence before this study}\par
\placeholder{We searched PubMed, Embase and arXiv/medRxiv from inception to
[date] for studies combining (multi-task OR transfer OR pretrain* OR foundation
model OR shared representation) AND (prognos* OR risk prediction) AND (electronic
health record OR registry), and for prognostic indices in lymphoid cancers.
Report numbers screened and included.} Clinical prediction models are almost always built for one disease and one
outcome. Multi-task and transfer learning on health records have helped some
tasks and harmed others, often against weak baselines. Models pretrained on
population records predict many outcomes, but whether they improve specialised
prognosis over models trained within the specialty, or over tuned
gradient-boosted trees, is unknown. Lymphoid cancers rely on disease-specific
indices built from a few variables, and many uncommon diseases have none.\par
\textbf{Added value of this study}\par
To our knowledge, this is the first study to measure, within one specialty and
on the same patients, what prognosis gains from borrowing across diseases,
across outcomes and from the national population. Matched models separate each
source across 24 diseases, 69 outcomes and six horizons. One model learns how
much each prediction should borrow, and every task is compared with tuned
tabular models and every prognostic index in routine use. \placeholder{Principal test findings with effect sizes.}\par
\textbf{Implications of all the available evidence}\par
Biological boundaries between diseases need not limit prognostic learning.
Other diseases can supply the examples a rare disease lacks, and one model can
serve a specialty.
\placeholder{Implications for disease-specific indices and task-specific
models, and what the national population adds.} Better prediction is not
evidence that acting on it improves care.
}}

\operasection{Introduction}
Prognosis is central to patient care. In haematology, it guides the choice of
treatment, determines eligibility for clinical trials and shapes what patients
and their families can expect. For the lymphoid cancers, a group of blood
cancers including lymphomas, lymphocytic leukaemias and myeloma, prognostic
tools such as the IPI, NCCN-IPI and CLL-IPI each predict a single outcome in a
single disease from a few variables
\citep{IPI1993,zhou2014nccnipi,cllipi2016,palumbo2015riss}. This specificity
comes at a price: it restricts the signals, patients and outcomes such tools can
learn from. Learning opportunities are scarce, because many patient groups and
outcomes are rare in isolation and rarer still in combination
\citep{Riley2020}. Adding more patients with the same disease is seldom possible
when the disease itself is rare. To predict better, information must be borrowed
from elsewhere.

Even for a single disease and outcome, other diseases, outcomes and patient
histories can be informative. Haematological cancers are biologically distinct
yet share diagnostic work-up, treatments and complications, and frailty and
prior illness affect outcomes globally. Learning such patterns across patients
and tasks can improve prediction from few labelled examples
\citep{Baxter2000,Maurer2016,zhang2019metapred}. Borrowing can also harm.
In clinical prediction, training tasks jointly frequently degraded individual
tasks, even where multi-task pretraining helped \citep{mcdermott2020mtl}. The question is therefore not only whether to borrow,
but what to borrow and how, without losing clinically important differences.

The simplest way for a model to borrow is to learn the same outcome across
several diseases, or several outcomes within the same disease, so that each task
gains examples from the others. Such pooling can cost the idiosyncratic
differences between subgroups. Ideally, a model relies on borrowed information
where it must, as for rare diseases and outcomes, and on its own where that is
sufficient. Statistics strikes this balance by partial pooling
\citep{EfronMorris1975}, and neural models can learn it from the data
\citep{cui2024autodp,lee2022residualmoe}.

A more recent way to borrow is to pretrain on patient histories, so that a model
learns general patterns of how illness, investigation and treatment unfold
before it learns any outcome \citep{li2020behrt,rasmy2021med,steinberg2023motor}.
A model pretrained on a seemingly unrelated cohort can then pass much of this
learned structure to a downstream prognostic model. Pretrained on the records of
a whole population, it borrows from people far outside the specialty
\citep{savcisens2024life2vec,shmatko2025delphi}. Whether these approaches improve
on strong conventional alternatives, particularly when labels are scarce,
remains uncertain \citep{wornow2023shaky,worth2026state}.

Lymphoid cancers are a natural proving ground for borrowing information across
diseases and outcomes. Their many entities are individually uncommon
\citep{alaggio2022who,sant2010haemacare} yet share investigation and care. At the
start of treatment, clinicians weigh disease control alongside infection,
cytopenia, organ toxicity and admission, which existing indices address only in
part \citep{DLG2025DLBCL,Tilly2022POLARIX}. Testing borrowing requires all of
these diseases and outcomes under one roof, recorded from the same sources in the
same formats. The Danish setting makes this possible: DALY-CARE links nationwide
registers, health records and laboratory data for all Danish adults with
lymphoid cancers \citep{brieghel2025dalycare}. Few, if any, other resources do
so.

We studied 69 outcomes in 37,531 patients with 24 lymphoid cancers and asked
three questions, answered in the Results sections named
(\cref{fig:paradigms}).
\begin{enumerate}[leftmargin=1.5em,itemsep=2pt,topsep=3pt]
  \item \emph{Whether to borrow.} Do sequence models improve when they learn
  across diseases, outcomes or both, compared with one model per disease and
  outcome? How do they compare with XGBoost and the clinical indices?
  (\emph{Whether to borrow}, \emph{Practical alternatives} and
  \emph{Prognostic indices}.)
  \item \emph{What to borrow.} What do other diseases, other outcomes and the
  national population each add? (\emph{What to borrow}.)
  \item \emph{How to borrow.} Does letting each prediction balance its own
  corrections against borrowed ones improve on borrowing uniformly, and where
  does any gain come from? (\emph{How to borrow}.)
\end{enumerate}

\begin{figure}[H]
  \centering
  \resizebox{\textwidth}{!}{\begin{tikzpicture}[opera,x=1cm,y=1cm]
  \draw[divider] (5.55,.05)--(5.55,-6.18);
  \draw[divider] (11.15,.05)--(11.15,-6.18);

  % Shared semantic bands make the three paradigms directly comparable.
  \begin{scope}[on background layer]
    \draw[draw=OperaBorder,densely dashed,line width=.45pt] (-.42,-2.70)--(16.55,-2.70);
    \draw[draw=OperaBorder,densely dashed,line width=.45pt] (-.42,-4.72)--(16.55,-4.72);
  \end{scope}
  \node[eyebrow,anchor=east,text=OperaInkMuted] at (-.54,-2.15) {INPUT};
  \node[eyebrow,anchor=east,text=OperaInkMuted] at (-.54,-3.68) {MODEL};
  \node[eyebrow,anchor=east,text=OperaInkMuted] at (-.54,-5.62) {OUTPUT};

  % A — clinical score
  \node[panelletter,anchor=north west] at (0,0) {A};
  \node[eyebrow,anchor=north west] at (.48,-.03) {CLINICAL SCORE};
  \node[heading,anchor=north west] at (.48,-.36) {A diagnostic snapshot};
  \node[labelsmall,anchor=north west,align=left,text width=4.65cm] at (.48,-.78)
    {Few variables, fixed combination, one disease};
  \foreach \name/\x/\lab in {age/.72/Age,stage/1.68/Stage,
      ldh/2.65/LDH,ps/3.62/PS,sites/4.58/Sites}
    \node[circle,draw=OperaSienna!55,fill=OperaSienna!7,minimum size=7.4mm,
      font=\sffamily\scriptsize] (\name) at (\x,-2.15) {\lab};
  \node[stagewarm,minimum width=2.15cm,minimum height=7mm] (score) at (2.65,-3.70) {risk score};
  \foreach \name in {age,stage,ldh,ps,sites}
    \draw[-{Latex[length=1.5mm]},draw=OperaSienna,line width=.65pt]
      (\name.south)--(score.north);
  \draw[flowwarm,shorten >=1pt] (score.south)--(2.65,-5.25);
  % Clinical scores define discrete strata, not a continuous probability scale.
  \node[labelsmall] at (1.38,-5.48) {low};
  \node[labelsmall] at (2.65,-5.48) {intermediate};
  \node[labelsmall] at (3.92,-5.48) {high};
  \fill[OperaSienna!16] (.75,-6.00) rectangle (2.02,-5.72);
  \fill[OperaSienna!43] (2.02,-6.00) rectangle (3.28,-5.72);
  \fill[OperaSienna!82] (3.28,-6.00) rectangle (4.55,-5.72);
  \draw[white,line width=.8pt] (2.02,-6.00)--(2.02,-5.72);
  \draw[white,line width=.8pt] (3.28,-6.00)--(3.28,-5.72);
  %\node[captionnote,text width=4.4cm] at (2.65,-6.72)
  %  {Interpretable and portable, but most prior history is compressed away.};

  % B — tabular model
  \node[panelletter,anchor=north west] at (5.86,0) {B};
  \node[eyebrow,anchor=north west] at (6.34,-.03) {TABULAR MACHINE LEARNING};
  \node[heading,anchor=north west] at (6.34,-.36) {History becomes features};
  \node[labelsmall,anchor=north west,align=left,text width=4.55cm] at (6.34,-.78)
    {Longitudinal history, engineered summaries, one task};
  \draw[OperaInkMuted,line width=.55pt] (6.35,-2.15)--(10.38,-2.15);
  \foreach \x/\col in {6.55/OperaPeriwinkle,7.22/OperaTeal,7.95/OperaTerracotta,
      8.63/OperaTeal,9.42/OperaOchre,10.08/OperaPeriwinkle}
    \fill[\col] (\x,-2.15) circle (2.1pt);
  % One history, partitioned into look-back intervals for aggregation.
  \foreach \x in {6.55,8.42,9.56,10.18}
    \draw[draw=OperaTeal!68,line width=.65pt] (\x,-2.03)--(\x,-2.27);
  \node[font=\sffamily\tiny,text=OperaTeal] at (7.48,-2.43) {365 d};
  \node[font=\sffamily\tiny,text=OperaTeal] at (8.99,-2.43) {180 d};
  \node[font=\sffamily\tiny,text=OperaTeal] at (9.87,-2.43) {30 d};
  \node[stage,draw=OperaTeal!65,fill=OperaTeal!7,minimum width=2.80cm,
    minimum height=7mm] (ml) at (8.36,-3.70)
    {ML algorithm\\{\scriptsize e.g. XGBoost}};
  \draw[flowteal,shorten >=1pt] (8.36,-2.42)--(ml.north);
  \node[labelsmall] (prob) at (8.36,-5.48) {$P(\mathrm{outcome}_n\ \mathrm{by}\ t)$};
  \draw[flowteal,shorten >=1pt] (ml.south)--(prob.north);
  \foreach \i in {0,...,23}{
    \pgfmathsetmacro{\xa}{6.80+0.13*\i}
    \pgfmathsetmacro{\xb}{6.80+0.13*(\i+1)}
    \pgfmathtruncatemacro{\tone}{6+3*\i}
    \fill[OperaTeal!\tone] (\xa,-6.00) rectangle (\xb,-5.72);
  }
  %\node[captionnote,text width=4.4cm] at (8.36,-6.72)
  %  {A separate predictor is typically fitted for each outcome and horizon.};

  % C — sequence model
  \node[panelletter,anchor=north west] at (11.46,0) {C};
  \node[eyebrow,anchor=north west] at (11.94,-.03) {LONGITUDINAL SEQUENCE MODEL};
  \node[heading,anchor=north west] at (11.94,-.36) {The timeline remains visible};
  \node[labelsmall,anchor=north west,align=left,text width=4.75cm] at (11.94,-.78)
    {Ordered events, contextual states, reusable learning};
  \draw[OperaIndigo,line width=.65pt] (11.95,-2.15)--(16.25,-2.15);
  \foreach \x/\col/\lab in {12.12/OperaPeriwinkle/D,12.82/OperaTeal/L,
      13.52/OperaTerracotta/M,14.22/OperaOchre/V,14.92/OperaTeal/L,15.62/OperaLavender/P}
    \node[circle,draw=\col,fill=white,minimum size=5.8mm,font=\sffamily\tiny] at (\x,-2.15) {\lab};
  \draw[flowblue] (14.10,-2.51)--(14.10,-2.80);
  \node[stageblue,minimum width=3.00cm,minimum height=7mm] (enc) at (14.10,-3.23)
    {transformer encoder};
  \node[softpanel,draw=OperaIndigo!42,fill=OperaIndigo!4,
    minimum width=2.25cm,minimum height=5.5mm,font=\sffamily\tiny]
    (z) at (14.10,-4.20) {shared representation $z$};
  \draw[flowblue,shorten >=1pt] (enc.south)--(z.north);
  \node[labelsmall] (c1) at (12.20,-5.48) {$P($tox. by $t)$};
  \node[labelsmall] (c2) at (13.70,-5.48) {$P($anemia by $t)$};
  \node[font=\sffamily\scriptsize,text=OperaInkMuted] at (14.80,-5.48) {$\cdots$};
  \node[labelsmall] (cn) at (16.02,-5.48) {$P($outcome$_n$ by $t)$};
  \draw[flowblue,shorten >=1pt] (z)--(c1.north);
  \draw[flowblue,shorten >=1pt] (z)--(c2.north);
  \draw[flowblue,shorten >=1pt] (z)--(cn.north);
  % Each task emits its own continuous probability, aligned with panel B.
  \begin{scope}[shift={(11.72,0)}]
    \foreach \i in {0,...,11}{
      \pgfmathsetmacro{\xa}{0.08*\i}
      \pgfmathsetmacro{\xb}{0.08*(\i+1)}
      \pgfmathtruncatemacro{\tone}{6+7*\i}
      \fill[OperaIndigo!\tone] (\xa,-6.00) rectangle (\xb,-5.72);
    }
  \end{scope}
  \begin{scope}[shift={(13.22,0)}]
    \foreach \i in {0,...,11}{
      \pgfmathsetmacro{\xa}{0.08*\i}
      \pgfmathsetmacro{\xb}{0.08*(\i+1)}
      \pgfmathtruncatemacro{\tone}{6+7*\i}
      \fill[OperaIndigo!\tone] (\xa,-6.00) rectangle (\xb,-5.72);
    }
  \end{scope}
  \begin{scope}[shift={(15.54,0)}]
    \foreach \i in {0,...,11}{
      \pgfmathsetmacro{\xa}{0.08*\i}
      \pgfmathsetmacro{\xb}{0.08*(\i+1)}
      \pgfmathtruncatemacro{\tone}{6+7*\i}
      \fill[OperaIndigo!\tone] (\xa,-6.00) rectangle (\xb,-5.72);
    }
  \end{scope}
  %\node[captionnote,text width=4.55cm] at (14.10,-6.72)
  %  {The encoder may be pretrained; one representation can support many outcomes and horizons.};
\end{tikzpicture}
}\par\vspace{3mm}
  \resizebox{\textwidth}{!}{\begin{tikzpicture}[opera,x=1cm,y=1cm]
  \draw[divider] (7.25,.05)--(7.25,-9.45);

  \node[panelletter,anchor=north west] at (0,0) {D};
  \node[eyebrow,anchor=north west] at (.48,-.03) {SOURCES OF PREDICTIVE INFORMATION};
  \node[heading,anchor=north west] at (.48,-.38) {More information for one prediction problem};
  \node[labelsmall,anchor=north west,align=left,text width=5.9cm] at (.48,-.82)
    {Moving right adds information to the learning problem.};

  % Concrete before-to-after examples are easier to read than abstract axes.
  \node[eyebrow,text=OperaIndigo,anchor=east] at (1.58,-1.90) {PATIENTS};
  \foreach \x/\y in {1.92/-1.76,2.12/-2.03,2.35/-1.80}
    \fill[OperaIndigo!55] (\x,\y) circle (1.8pt);
  \draw[flowblue] (2.62,-1.90)--(3.42,-1.90);
  \foreach \x/\y in {3.72/-1.70,3.93/-2.02,4.15/-1.78,4.38/-2.07,4.58/-1.69,
      4.82/-1.94,5.05/-1.74,5.28/-2.06,5.52/-1.84}
    \fill[OperaIndigo!65] (\x,\y) circle (1.8pt);
  \node[labelsmall,anchor=west] at (1.82,-2.27) {larger sample};

  \node[eyebrow,text=OperaTeal,anchor=east] at (1.58,-3.15) {VARIABLES};
  \draw[OperaTeal!55,line width=.55pt] (1.86,-3.41) rectangle (2.52,-2.89);
  \draw[OperaTeal!55,line width=.45pt] (2.18,-3.41)--(2.18,-2.89);
  \draw[flowteal] (2.72,-3.15)--(3.42,-3.15);
  \draw[OperaTeal!65,line width=.55pt] (3.70,-3.51) rectangle (5.48,-2.79);
  \foreach \x in {4.05,4.40,4.75,5.10}\draw[OperaTeal!45,line width=.4pt] (\x,-3.51)--(\x,-2.79);
  \node[labelsmall,anchor=west] at (1.82,-3.72) {richer phenotype};

  \node[eyebrow,text=OperaSienna,anchor=east] at (1.58,-4.45) {TIME};
  \fill[OperaSienna!65] (2.18,-4.45) circle (2.2pt);
  \draw[flowwarm] (2.72,-4.45)--(3.42,-4.45);
  \draw[OperaSienna!70,line width=.65pt] (3.70,-4.45)--(5.48,-4.45);
  \foreach \x in {3.82,4.18,4.55,4.94,5.34}\fill[OperaSienna!70] (\x,-4.45) circle (1.9pt);
  \node[labelsmall,anchor=west] at (1.82,-4.85) {longer history};

  \draw[draw=OperaBorder,densely dashed,line width=.5pt] (1.82,-5.16)--(5.62,-5.16);

  \node[eyebrow,text=OperaIndigo,anchor=east] at (1.58,-5.72) {DISEASES};
  \foreach \x/\y/\col in {1.98/-5.63/OperaIndigo,2.28/-5.81/OperaPeriwinkle}
    \fill[\col!70] (\x,\y) circle (2.3pt);
  \draw[flowblue] (2.72,-5.72)--(3.42,-5.72);
  \foreach \x/\y/\col in {3.72/-5.58/OperaIndigo,4.05/-5.82/OperaPeriwinkle,
      4.38/-5.62/OperaTeal,4.72/-5.84/OperaLavender,5.05/-5.57/OperaIndigo,
      5.38/-5.79/OperaTeal,5.66/-5.62/OperaPeriwinkle}
    \fill[\col!72] (\x,\y) circle (2.3pt);
  \node[labelsmall,anchor=west] at (1.82,-6.08) {broader cohort space};

  \node[eyebrow,text=OperaTeal,anchor=east] at (1.58,-6.82) {OUTCOMES};
  \foreach \x/\y in {2.02/-6.72,2.32/-6.92}
    \fill[OperaTeal!55] (\x-.08,\y-.08) rectangle (\x+.08,\y+.08);
  \draw[flowteal] (2.72,-6.82)--(3.42,-6.82);
  \foreach \x/\y/\tone in {3.72/-6.68/38,4.05/-6.92/48,4.38/-6.72/58,
      4.72/-6.94/68,5.05/-6.67/44,5.38/-6.88/76,5.66/-6.70/56}
    \fill[OperaTeal!\tone] (\x-.08,\y-.08) rectangle (\x+.08,\y+.08);
  \node[labelsmall,anchor=west] at (1.82,-7.18) {more learning tasks};
  \node[eyebrow,text=OperaGreen,anchor=east] at (1.58,-7.48) {MORE INFORMATION};
  \draw[-{Latex[length=2.2mm]},draw=OperaGreen,line width=1pt] (1.82,-7.48)--(5.48,-7.48);
  \node[labelsmall,text=OperaGreen] at (3.65,-7.76)
    {higher attainable performance};

  \node[panelletter,anchor=north west] at (7.60,0) {E};
  \node[eyebrow,anchor=north west] at (8.08,-.03) {LEARNING ACROSS COHORTS AND TASKS};
  \node[heading,anchor=north west] at (8.08,-.38) {Related diseases and outcomes};
  \node[labelsmall,anchor=north west,align=left,text width=7.6cm] at (8.08,-.82)
    {More disease cohorts and more outcome tasks add structured information around a low-support target.};

  % Separate the axis/title and matrix from the two-line panel introduction.
  \begin{scope}[yshift=-5mm]
  % Keep the axis title separate from the column labels and first row.
  \draw[-{Latex[length=1.8mm]},draw=OperaIndigo,line width=.75pt]
    (10.36,-1.55)--(16.72,-1.55);
  \node[eyebrow,text=OperaIndigo] at (13.54,-1.30) {MORE OUTCOME TASKS};

  \foreach \r/\disease in {0/DLBCL,1/Multiple myeloma,2/CLL,3/Hodgkin lymphoma,4/PTCL}
    \node[labelsmall,anchor=east] at (10.20,-2.72-.84*\r) {\disease};
  \draw[-{Latex[length=1.8mm]},draw=OperaTeal,line width=.75pt]
    (7.76,-2.40)--(7.76,-7.62);
  \node[eyebrow,text=OperaTeal,rotate=90] at (7.47,-5.01) {MORE DISEASE COHORTS};
  \foreach \col/\outcome in {0/Anaemia,1/Infection,2/Renal,3/Fracture,4/Survival}
    \node[font=\sffamily\fontsize{6.2}{7}\selectfont,text=OperaInkSecondary,
      align=center,text width=1.05cm,anchor=south] at (10.82+1.12*\col,-2.27) {\outcome};
  \node[font=\sffamily\fontsize{6.2}{7}\selectfont,text=OperaInkSecondary,
    align=center,text width=1.05cm,anchor=south] at (16.72,-2.27) {Outcome$_n$};
  \foreach \r in {0,...,4}{\foreach \col in {0,...,4}{
    \pgfmathtruncatemacro{\tone}{7+3*\r+2*\col}
    \node[rounded corners=.35mm,draw=OperaBorder,fill=OperaIndigo!\tone,
      minimum width=8.8mm,minimum height=5.7mm] (m\r\col)
      at (10.82+1.12*\col,-2.72-.84*\r) {};
  }}
  \foreach \r in {0,...,4}{
    \node[font=\sffamily\scriptsize,text=OperaInkMuted]
      at (16.01,-2.72-.84*\r) {$\cdots$};
  }
  \foreach \r in {0,...,4}{
    \pgfmathtruncatemacro{\tone}{16+3*\r}
    \node[rounded corners=.35mm,draw=OperaBorder,fill=OperaIndigo!\tone,
      minimum width=8.8mm,minimum height=5.7mm] (mn\r) at (16.72,-2.72-.84*\r) {};
  }
  \foreach \x in {10.82,11.94,13.06,14.18,15.30,16.72}{
    \node[font=\sffamily\scriptsize,text=OperaInkMuted] at (\x,-6.60) {$\vdots$};
  }
  \node[font=\sffamily\scriptsize,text=OperaInkMuted] at (16.01,-6.60) {$\ddots$};
  \node[labelsmall,anchor=east] at (10.20,-7.29) {Disease$_m$};
  \foreach \col in {0,...,4}{
    \pgfmathtruncatemacro{\tone}{22+2*\col}
    \node[rounded corners=.35mm,draw=OperaBorder,fill=OperaIndigo!\tone,
      minimum width=8.8mm,minimum height=5.7mm] (mm\col)
      at (10.82+1.12*\col,-7.29) {};
  }
  \node[rounded corners=.35mm,draw=OperaBorder,fill=OperaIndigo!34,
    minimum width=8.8mm,minimum height=5.7mm] (mmn) at (16.72,-7.29) {};
  \node[rounded corners=.35mm,draw=OperaSienna,line width=1.1pt,fill=OperaSienna!8,
    minimum width=8.8mm,minimum height=5.7mm] (target) at (11.94,-6.08) {};
  % Every training cell is connected; most links recede behind the emphasized
  % same-outcome and same-disease borrowing paths.
  \begin{scope}[on background layer]
    \foreach \r in {0,...,4}{\foreach \col in {0,...,4}{
      \draw[draw=OperaTeal,opacity=.55,line width=.6pt] (m\r\col)--(target);
    }}
    \foreach \r in {0,...,4}{
      \draw[draw=OperaTeal,opacity=.55,line width=.6pt] (mn\r)--(target);
    }
    \foreach \col in {0,...,4}{
      \draw[draw=OperaTeal,opacity=.55,line width=.6pt] (mm\col)--(target);
    }
    \draw[draw=OperaTeal,opacity=.55,line width=.6pt] (mmn)--(target);
    \foreach \r in {0,...,3}{
      \draw[draw=OperaTeal!85!black,opacity=1,line width=1.6pt] (m\r1)--(target);
    }
    \foreach \col in {0,2,3,4}{
      \draw[draw=OperaTeal!85!black,opacity=1,line width=1.6pt] (m4\col)--(target);
    }
    % The last disease (same outcome) and last outcome (same disease) borrow too.
    \draw[draw=OperaTeal!85!black,opacity=1,line width=1.6pt] (mm1)--(target);
    \draw[draw=OperaTeal!85!black,opacity=1,line width=1.6pt] (mn4)--(target);
  \end{scope}
  \end{scope}
\end{tikzpicture}
}
  \caption{\textbf{Approaches to risk prediction and the information each
  can use.} Top, three approaches to risk prediction in lymphoid cancers.
  \textbf{A}, clinical risk scores combine a small diagnostic snapshot with a
  fixed rule into broad risk strata. \textbf{B}, classical machine learning
  uses a wider history, aggregated over time into a task-specific feature
  table. \textbf{C}, a longitudinal sequence model reads the ordered clinical
  events and can share one patient representation across outcomes. Its
  encoder may be pretrained on the cohort or on the whole population. Bottom,
  sources of additional signal. \textbf{D}, conventional scaling adds patients,
  variables or observations to one isolated prediction problem. \textbf{E},
  shared learning also draws on related diseases (rows) and outcomes
  (columns).}
  \label{fig:paradigms}
\end{figure}

\clearpage
\operasection{Methods}

\operasubsection{Population, prediction time and outcomes}
We included all patients in Denmark diagnosed with a lymphoid cancer
(lymphoma, lymphocytic leukaemia or plasma-cell neoplasm) from 2002 who were
aged 18 or older at diagnosis and started first-line treatment, as recorded in
DALY-CARE \citep{brieghel2025dalycare}. Every prediction was made at the start of
treatment, based only on records from before it. These comprised diagnoses,
procedures, pathology, medication, laboratory results and hospital contacts with
their region and specialty of care, together with age and sex
(\cref{fig:model-inputs}). The study was conducted under the DALY-CARE approvals
from the Danish Health Data Authority (P-2020-561) and the National Ethics
Committee (1804410).

We predicted 69 outcomes at six horizons from 30 days to two years, spanning
disease course, infection, cytopenia, organ toxicity, cardiovascular events and
acute care, with death as a competing event for non-fatal outcomes. Laboratory
outcomes followed the Common Terminology Criteria for Adverse Events (CTCAE)
\citep{nci2017ctcae} as closely as routine data permitted. These are graded
abnormal laboratory values and not necessarily toxicities caused by treatment.
Definitions, adaptations and ascertainment, including which patients were
eligible for each laboratory outcome, are given in the Supplementary Methods.

We used a pseudo-prospective design, splitting patients by the date treatment
started. Models were trained on patients treated from 2002 to 2021, tuned on those
treated in 2022 (the tuning set) and tested on those treated from 2023, a set
left unopened until every choice was fixed. Pretraining used only records from
before 2022.
\previewonly{The current figures and numbers are tuning analyses. Held-out and
national-pretraining results remain pending.}

\operasubsection{Borrowing across diseases and outcomes}
To investigate the effect of borrowing across diseases and outcomes, we
organised all 1,656 prediction tasks (24 diseases by 69 outcomes) in a grid
(\cref{fig:supervision-grid}). Each cell is a prediction question, such as the
risk of infection in a patient with DLBCL, at six horizons. Four matched types of model differed only in what they could
borrow. A specialist learned from one cell. A disease-pooled model learned one
outcome across all diseases. An outcome-pooled model learned all outcomes within
one disease. A joint model learned across both diseases and outcomes. All shared a causal transformer encoder
pretrained to predict the next recorded event, building on CORE-BEHRT
\citep{odgaard2024core}, and were fine-tuned to estimate cumulative incidence at
each horizon, with death as a competing event
\citep{gensheimer2019nnetsurvival,lee2018deephit}. Every outcome was predicted
and evaluated separately (Supplementary Methods).

\operasubsection{Learning how much to borrow}
Borrowing uniformly can blur differences that matter, by assuming that diseases
and outcomes are similar when they are not. Outcome Prediction with Expert Routing Across diseases and outcomes
(\opera{}) instead lets each prediction balance shared and task-specific
information (\cref{fig:source-router}). We added small trainable modules, called
adapters, that correct the shared patient representation: one for each disease,
one for each outcome, and one shared by all tasks. A router weighs the disease
and outcome corrections for each patient and target outcome, so that a
prediction can rely on its own where they are informative and borrow where its
own data are scarce \citep{ma2018mmoe,houlsby2019adapters}. The idea is that of partial pooling
\citep{EfronMorris1975}, with the balance learned from the data rather than set
by a prior. Model selection is described in the Supplementary Methods.

To understand how these corrections helped predictions, we constructed three
separately trained controls. The first keeps only the shared correction, the
second adds the patient's own disease and outcome corrections, and the third,
the full model, also borrows those of other tasks. Replacing other diseases' corrections
by their averages showed how much the fitted model depended on each
(Supplementary Methods).

\operasubsection{Borrowing from the population}
To test the limits of borrowing, we pretrained a causal transformer, PHAIR, not
only on patients with lymphoid cancers but on the nationwide Danish registers,
covering around \placeholder{8--9 million} people, using records from before
1 January 2022. We fine-tuned the joint model and
\opera{} from it directly and after continued pretraining on the cohort's records
\citep{gururangan2020dont}. The encoder used by \opera{} in all other
comparisons was chosen among these on the tuning set (Supplementary Methods).
The encoders share registers, architecture and training schedule, but the
national vocabulary keeps rarer codes and finer laboratory bins, so the
comparison concerns pretraining strategies rather than the number of patients
alone (Supplementary Methods). Because pretraining should matter most when labels
are scarce, joint models were also fine-tuned on 10\%, 25\% and 50\% of the
training patients, mimicking smaller cohorts.

\operasubsection{Comparators}
We compared the sequence models with XGBoost \citep{chen2016xgboost} and ridge
logistic regression, which remain hard to outperform on routine records
\citep{grinsztajn2022tabular}. Features summarised the same histories over
\placeholder{7-, 30-, 90- and 365-day} windows, using code counts and the latest,
minimum, maximum, variance and slope of laboratory values
\citep{bernstorff2023tsflattener}. This aggregation already improved on the
clinical prognostic indices in DLBCL \citep{werling2026dlbcl}. A separate model
was tuned for each disease, outcome and horizon. We also made versions of the
XGBoost models that pooled across diseases.

Each clinical index was compared with the model on the same patients. Ordered
by predicted risk, the model placed patients into as many groups as the index
has categories, each the size of the corresponding category, and survival was
compared between the two sets of groups. No outcome information set these
boundaries. These comparisons evaluate the model as a whole, including its use
of the full record, rather than borrowing alone.

\operasubsection{Evaluation}
We assessed discrimination with the time-dependent area under the ROC curve
(AUC) at each horizon and probability accuracy with the Brier score and index of
prediction accuracy (IPA), with censoring weights
\citep{blanche2013competingauc,gerds2006brier,kattan2018ipa}. Calibration was
assessed at every horizon, without recalibrating any model. Comparisons were
paired on the same patients and restricted to disease--outcome--horizon cells
with at least five cases and five controls, the minimum needed to estimate an
AUC. Cells meeting it can still be imprecise. Because diseases differ greatly in
size, and larger diseases contribute more cells, every summary is reported two
ways:
with each supported cell weighing equally, and with each disease weighing
equally (endpoint families, outcomes and horizons balanced within it). Intervals
came from 1,000 patient bootstrap resamples within disease, shared across
models (Supplementary Methods). Integrated gradients, and evaluations with
recent weeks or entire sources of the record hidden, described what the
predictions depend on. Reporting follows
TRIPOD+AI \citep{collins2024tripodai}.

\operasubsection{Role of the funding source}
\placeholder{The funders' roles in study design, data collection, analysis,
interpretation, writing and the decision to submit.}

\FloatBarrier
\operasection{Results}

\operasubsection{Patients and the tasks that could be evaluated}
Of 78,668 adults in DALY-CARE, whose records before 2022 pretrained the cohort
encoder, 37,531 started first-line treatment with an assigned disease
(\cref{fig:cohort-flow,tab:characteristics}): 29,105 were treated through 2021,
1,980 in 2022 and 6,446 from 2023. The median age was 69 years, and diffuse
large B-cell lymphoma and multiple myeloma made up half of the cohort.
Laboratory records and the third version of the National Patient Register
(LPR3, from 2019) came online late in the training period, so later patients had
longer, more laboratory-rich records. They were also older and more often lacked
a recorded clinical score, so the test evaluated performance under a shift in
both patients and records. \result{P5_cells}{N} disease--outcome--horizon
cells, \result{P5_cells_1y}{N} of them at one year, had enough events for the
specialist and joint models to be compared
(\cref{fig:supp-coverage,fig:supp-support}). % Only the patient-count row below has been updated for the released supervised population.
% All other numeric rows remain historical layout values and MUST be re-exported before use.
\begin{table*}[t]
\centering
\caption{Baseline characteristics of the supervised first-line haematology analysis cohort. \placeholder{All characteristics except patient counts require re-export from the current release.}}
\label{tab:characteristics}
\footnotesize
\setlength{\tabcolsep}{4pt}
\begin{tabularx}{\textwidth}{@{}Xrrrr@{}}
\toprule
Characteristic & Overall & Training & Tuning & Held-out \\
\midrule
Patients, $n$ & 37,531 & 29,105 & 1,980 & 6,446 \\
Age, years, median (IQR) & 69.3 (59.7--77.1) & 68.7 (59.2--76.5) & 71.1 (61.1--78.4) & 71.7 (61.7--78.8) \\
Age $\geq$70 years, $n$ (\%) & 18,011 (48.1\%) & 13,431 (46.1\%) & 1,050 (53.4\%) & 3,530 (55.2\%) \\
Female, $n$ (\%) & 15,647 (41.8\%) & 12,199 (41.9\%) & 802 (40.8\%) & 2,646 (41.4\%) \\
Male, $n$ (\%) & 21,820 (58.2\%) & 16,909 (58.1\%) & 1,165 (59.2\%) & 3,746 (58.6\%) \\
Year of first-line treatment, median (IQR) & 2016 (2010--2021) & 2014 (2009--2018) & 2022 (2022--2022) & 2024 (2023--2025) \\
Observed follow-up, years, median (IQR) & 4.1 (1.4--8.6) & 5.8 (2.1--10.1) & 3.5 (3.2--3.8) & 1.3 (0.6--2.2) \\
\addlinespace
\multicolumn{5}{@{}l}{\sffamily\bfseries Disease groups, $n$ (\%)}\\
Amyloidosis & 47 (0.1\%) & 22 (0.1\%) & 7 (0.4\%) & 18 (0.3\%) \\
\textbf{Burkitt/lymphoblastic lymphoma} & 1,364 (3.6\%) & 1,075 (3.7\%) & 87 (4.4\%) & 202 (3.2\%) \\
\hspace*{1em}Burkitt lymphoma (BL) & 295 (0.8\%) & 238 (0.8\%) & 15 (0.8\%) & 42 (0.7\%) \\
\hspace*{1em}Lymphoblastic lymphoma (LBL) & 1,069 (2.9\%) & 837 (2.9\%) & 72 (3.7\%) & 160 (2.5\%) \\
\textbf{CLL/SLL} & 5,152 (13.8\%) & 4,092 (14.1\%) & 244 (12.4\%) & 816 (12.8\%) \\
\hspace*{1em}Chronic lymphocytic leukaemia (CLL) & 4,062 (10.8\%) & 3,220 (11.1\%) & 187 (9.5\%) & 655 (10.2\%) \\
\hspace*{1em}Small lymphocytic lymphoma (SLL) & 1,090 (2.9\%) & 872 (3.0\%) & 57 (2.9\%) & 161 (2.5\%) \\
\textbf{DLBCL and related aggressive B-cell lymphomas} & 10,408 (27.8\%) & 8,369 (28.8\%) & 505 (25.7\%) & 1,534 (24.0\%) \\
\hspace*{1em}Diffuse large B-cell lymphoma (DLBCL) & 9,951 (26.6\%) & 8,008 (27.5\%) & 481 (24.5\%) & 1,462 (22.9\%) \\
\hspace*{1em}Richter transformation (RT) & 192 (0.5\%) & 151 (0.5\%) & 14 (0.7\%) & 27 (0.4\%) \\
\hspace*{1em}Transformed follicular lymphoma & 265 (0.7\%) & 210 (0.7\%) & 10 (0.5\%) & 45 (0.7\%) \\
Hairy-cell leukaemia & 418 (1.1\%) & 293 (1.0\%) & 29 (1.5\%) & 96 (1.5\%) \\
Hodgkin lymphoma & 2,644 (7.1\%) & 2,080 (7.1\%) & 141 (7.2\%) & 423 (6.6\%) \\
\textbf{Indolent B-cell non-Hodgkin lymphoma} & 5,399 (14.4\%) & 3,895 (13.4\%) & 302 (15.4\%) & 1,202 (18.8\%) \\
\hspace*{1em}Follicular lymphoma (FL) & 2,932 (7.8\%) & 2,248 (7.7\%) & 150 (7.6\%) & 534 (8.4\%) \\
\hspace*{1em}B-cell lymphoma, unspecified (BCL) & 227 (0.6\%) & 87 (0.3\%) & 17 (0.9\%) & 123 (1.9\%) \\
\hspace*{1em}Lymphoplasmacytic lymphoma (LPL) & 1,570 (4.2\%) & 1,165 (4.0\%) & 100 (5.1\%) & 305 (4.8\%) \\
\hspace*{1em}Extranodal marginal-zone lymphoma (EMZL) & 500 (1.3\%) & 318 (1.1\%) & 33 (1.7\%) & 149 (2.3\%) \\
\hspace*{1em}Nodal marginal-zone lymphoma (NMZL) & 27 (0.1\%) & 5 (0.0\%) & $<$5 & 22 (0.3\%) \\
\hspace*{1em}Splenic marginal-zone lymphoma (SMZL) & 143 (0.4\%) & 72 (0.2\%) & $<$5 & 69 (1.1\%) \\
Mantle-cell lymphoma & 1,466 (3.9\%) & 1,116 (3.8\%) & 87 (4.4\%) & 263 (4.1\%) \\
\textbf{Myeloma and related plasma-cell disorders} & 8,953 (23.9\%) & 6,904 (23.7\%) & 491 (25.0\%) & 1,558 (24.4\%) \\
\hspace*{1em}Multiple myeloma (MM) & 8,629 (23.0\%) & 6,673 (22.9\%) & 464 (23.6\%) & 1,492 (23.3\%) \\
\hspace*{1em}Plasma-cell leukaemia (PCL) & 201 (0.5\%) & 146 (0.5\%) & 18 (0.9\%) & 37 (0.6\%) \\
\hspace*{1em}Solitary myeloma/plasmacytoma & 123 (0.3\%) & 85 (0.3\%) & 9 (0.5\%) & 29 (0.5\%) \\
\textbf{T-cell non-Hodgkin lymphoma} & 1,616 (4.3\%) & 1,262 (4.3\%) & 74 (3.8\%) & 280 (4.4\%) \\
\hspace*{1em}Angioimmunoblastic T-cell lymphoma (AITL) & 483 (1.3\%) & 351 (1.2\%) & 26 (1.3\%) & 106 (1.7\%) \\
\hspace*{1em}Anaplastic large-cell lymphoma (ALCL) & 410 (1.1\%) & 339 (1.2\%) & 12 (0.6\%) & 59 (0.9\%) \\
\hspace*{1em}Peripheral T-cell lymphoma (PTCL) & 522 (1.4\%) & 444 (1.5\%) & 17 (0.9\%) & 61 (1.0\%) \\
\hspace*{1em}T-cell lymphoma, unspecified (TCL) & 201 (0.5\%) & 128 (0.4\%) & 19 (1.0\%) & 54 (0.8\%) \\
\addlinespace
\multicolumn{5}{@{}l}{\sffamily\bfseries Recorded performance status, $n$ (\%)}\\
Performance status 0 & 14,662 (53.8\%) & 11,842 (53.2\%) & 858 (54.3\%) & 1,962 (57.9\%) \\
Performance status 1 & 8,262 (30.3\%) & 6,838 (30.7\%) & 434 (27.5\%) & 990 (29.2\%) \\
Performance status 2 & 2,557 (9.4\%) & 2,077 (9.3\%) & 187 (11.8\%) & 293 (8.6\%) \\
Performance status 3 & 1,273 (4.7\%) & 1,063 (4.8\%) & 88 (5.6\%) & 122 (3.6\%) \\
Performance status 4 & 481 (1.8\%) & 444 (2.0\%) & 13 (0.8\%) & 24 (0.7\%) \\
Missing performance status$^{a}$ & 10,232 (27.3\%) & 6,844 (23.5\%) & 387 (19.7\%) & 3,001 (46.9\%) \\
\addlinespace
\multicolumn{5}{@{}l}{\sffamily\bfseries Disease-specific prognostic risk category, $n$ (\%)}\\
Low & 3,989 (17.6\%) & 3,231 (17.4\%) & 243 (17.6\%) & 515 (18.9\%) \\
Intermediate & 10,513 (46.4\%) & 8,608 (46.4\%) & 639 (46.3\%) & 1,266 (46.5\%) \\
High / very high$^{b}$ & 8,161 (36.0\%) & 6,721 (36.2\%) & 497 (36.0\%) & 943 (34.6\%) \\
Missing risk category$^{a}$ & 14,804 (39.5\%) & 10,548 (36.2\%) & 588 (29.9\%) & 3,668 (57.4\%) \\
\addlinespace
\multicolumn{5}{@{}l}{\sffamily\bfseries Pre-treatment record summaries, median (IQR)$^{c}$}\\
Recorded clinical-history span, years & \placeholder{Re-export} & \placeholder{Re-export} & \placeholder{Re-export} & \placeholder{Re-export} \\
Clinical-coded records & \placeholder{Re-export} & \placeholder{Re-export} & \placeholder{Re-export} & \placeholder{Re-export} \\
Distinct diagnosis codes & 21.0 (13.0--32.0) & 20.0 (12.0--30.0) & 24.0 (15.0--37.0) & 25.0 (15.0--38.0) \\
Distinct medication codes & 21.0 (12.0--34.0) & 18.0 (10.0--30.0) & 32.0 (21.0--46.0) & 34.0 (22.0--49.0) \\
Laboratory-coded records & 197.0 (0.0--482.0) & 118.0 (0.0--346.0) & 483.0 (279.5--803.5) & 531.0 (302.0--887.0) \\
\bottomrule
\end{tabularx}
\vspace{1.5mm}
\begin{minipage}{\textwidth}\scriptsize
One record per patient in the historical 37,467-person review export; these values are not the current 37,557-person first-line source or its supervised analysis subset. Disease groups are from \texttt{cohort\_grouped}, not the separate, incomplete \texttt{Disease} field. Indented rows show fine cohorts within their grouped parent; one-to-one groups are shown once. Disease percentages use the full cohort/split, not the parent group. A dash denotes an absent row in the source export, not a confirmed zero. Detailed treatment categories and scoring systems are reported in the supplementary tables.
$^{a}$Performance status and disease-specific prognostic risk categories are sourced from quality registries, whose disease and calendar-period coverage is incomplete. Category percentages use patients with nonmissing values; missingness percentages use the full cohort/split. The later-period decline in availability must be re-estimated from the current release before interpretation.
$^{b}$High and very high are combined only for this descriptive overview. The very-high category occurs in CLL-IPI (114 patients overall); risk categories from different scoring systems are not equivalent clinical scales. Original categories and scoring-system identities are retained in the supplement.
$^{c}$MEDS-derived values are provisional pending re-export after excluding demographic tokens from clinical-history summaries. Laboratory entries count coded records, not unique tests or verified source measurements. Follow-up is observed duration, not a reverse-Kaplan--Meier estimate.
IQR, interquartile range. Primary-only small-cell masking in the source exports is not disclosure clearance; the complete publication bundle requires institutional review.
\end{minipage}
\end{table*}


\FloatBarrier
\operasubsection{Whether to borrow: borrowing against task-specific models}
Borrowing from other tasks improved the sequence models. Compared with
specialists, AUC changed by \result{P1_auc}{X} when models borrowed from other
diseases (disease-balanced \result{P1_auc_disease_balanced}{X}), by
\result{P2_auc}{Y} when they borrowed from other outcomes
(\result{P2_auc_disease_balanced}{Y}) and by \result{P5_auc}{Z} when they
borrowed from both (\result{P5_auc_disease_balanced}{Z}). IPA changed by
\result{P1_ipa}{A}, \result{P2_ipa}{B} and \result{P5_ipa}{C}
(\cref{fig:pooling}, \cref{tab:panel-effects}). The joint model had the higher
AUC in \result{P5_better}{N of M} cells (\resultpct{P5_better}). \placeholder{The domains and horizons in which the
gain was largest and smallest} (\cref{fig:pooling}D, \cref{fig:supp-horizons}).

The gain from borrowing was largest where a task had least data of its own
\placeholder{(confirm on test)} (\cref{fig:pooling}C). In diseases with fewer
than 500 training patients, the joint model differed from the specialist by
\result{P5_auc_rare}{X}. Borrowing also produced predictions where no
specialist could be trained. At one year, \result{cells_no_training_events}{102}
disease--outcome cells had no training events and \result{cells_fit_failed}{19}
specialist fits failed, but the joint model predicted all of them, and
\result{cells_joint_only_evaluable}{N} had enough events to be evaluated
(\cref{fig:supp-coverage,fig:supp-pooling-cells}).

\begin{figure}[H]
  \resultfigure{fig_pooling}{Pooling across diseases and outcomes}
    {A: joint versus specialist by disease and domain. B: gain by pooling axis.
    C: gain against training events. D: gain by clinical domain.}
  \caption{\textbf{Effect of pooling across diseases and outcomes on
  discrimination.} \textbf{A}, AUC gain of the joint model over the
  specialist, pooled over the six horizons, by disease (rows) and clinical
  domain (columns), median over the domain's outcomes. Diseases run from fewest to most training patients, given beside each name with a log-scale bar, and domains from fewest to most training events.
  \textbf{B}, For each disease, the mean gain from pooling
  over diseases, pooling over outcomes and both (joint), on the same rows.
  \textbf{C}, Gain against the number of training events of the target cell
  (at one year),
  with the hierarchical support curve and 95\% credible band. Points are cells
  with at least 10 training events and a 95\% interval narrower than 0.4 (all
  cells enter the curve). \previewonly{Tuning preview: the support curve is a kernel smoother over cells with a 95\% bootstrap band. The hierarchical model is fitted for the test-set analysis.} \textbf{D}, Gains
  by clinical domain (paired mean and 95\% bootstrap interval where
  available, with descriptive preview fallbacks labelled in the plot). Positive
  values favour the pooled model, and grey cells could not be estimated.
  Cell-level maps for each pooling axis, their
  interaction and all horizons are in
  \cref{fig:supp-pooling-cells,fig:supp-horizons}.}
  \label{fig:pooling}
\end{figure}

\begin{table}[H]
\centering\small
\caption{\textbf{What each source of signal adds, across all horizons.}
Paired differences on the test set (first model minus second), pooled over all
six horizons with equal weight per disease--outcome--horizon cell, with 95\%
paired bootstrap intervals. The same table at each horizon is in Supplementary
Table~\placeholder{S}. Cells:
cells estimable for both models. Rare diseases: fewer than 500 training
patients. Cells better: share of cells with a higher AUC point estimate,
irrespective of uncertainty. Bold: the interval lies above zero. Red: below
zero. Claim: the prespecified label from the paired bootstrap of the pooled
difference: \emph{better} (\emph{worse}) if the interval lies above (below)
zero, marked \emph{above margin} when the estimate also exceeds the practical
margin of 0.01. \emph{Equivalent} if the interval lies within $\pm0.01$,
otherwise \emph{inconclusive}. Paired bootstrap $p$-values and the hierarchical
estimates, with their separately specified weighting, are in Supplementary
Table~\placeholder{S}.}
\label{tab:panel-effects}
\resizebox{\textwidth}{!}{%
\begin{tabular}{@{}lrrrrrl@{}}
\toprule
Comparison & Cells & $\Delta$AUC (95\% CI) & $\Delta$IPA (95\% CI) & $\Delta$AUC, rare diseases & Cells better & Claim \\
\midrule
\IfFileExists{\resultnumbersdir panel_effects_rows.tex}{\input{\resultnumbersdir panel_effects_rows.tex}}{%
\rlap{\textit{Learning across diseases and outcomes (against the specialist)}} & & & & & & \\
\quad Other diseases (disease-pooled) & \placeholder{N} & \placeholder{x} & \placeholder{x} & \placeholder{x} & \placeholder{x\%} & \placeholder{} \\
\quad Other outcomes (outcome-pooled) & \placeholder{N} & \placeholder{x} & \placeholder{x} & \placeholder{x} & \placeholder{x\%} & \placeholder{} \\
\quad Both (joint) & \placeholder{N} & \placeholder{x} & \placeholder{x} & \placeholder{x} & \placeholder{x\%} & \placeholder{} \\
\rlap{\textit{Learning how much to borrow}} & & & & & & \\
\quad Routed vs joint & \placeholder{N} & \placeholder{x} & \placeholder{x} & \placeholder{x} & \placeholder{x\%} & \placeholder{} \\
\rlap{\textit{Population pretraining (joint model)}} & & & & & & \\
\quad National vs cohort encoder & \placeholder{N} & \placeholder{x} & \placeholder{x} & \placeholder{x} & \placeholder{x\%} & \placeholder{} \\
\quad Adapted national vs cohort encoder & \placeholder{N} & \placeholder{x} & \placeholder{x} & \placeholder{x} & \placeholder{x\%} & \placeholder{} \\
\rlap{\textit{Population pretraining (OPERA)}} & & & & & & \\
\quad \opera{} (national) vs \opera{} (local) & \placeholder{N} & \placeholder{x} & \placeholder{x} & \placeholder{x} & \placeholder{x\%} & \placeholder{} \\
\quad \opera{} (national + adaptation) vs \opera{} (local) & \placeholder{N} & \placeholder{x} & \placeholder{x} & \placeholder{x} & \placeholder{x\%} & \placeholder{} \\
\rlap{\textit{Sequence models against tabular models}} & & & & & & \\
\quad Specialist vs per-task XGBoost & \placeholder{N} & \placeholder{x} & \placeholder{x} & \placeholder{x} & \placeholder{x\%} & \placeholder{} \\
\quad Pooled vs per-task XGBoost & \placeholder{N} & \placeholder{x} & \placeholder{x} & \placeholder{x} & \placeholder{x\%} & \placeholder{} \\
\quad \opera{} (national + adaptation) vs per-task XGBoost & \placeholder{N} & \placeholder{x} & \placeholder{x} & \placeholder{x} & \placeholder{x\%} & \placeholder{} \\
\quad \opera{} (local) vs per-task XGBoost & \placeholder{N} & \placeholder{x} & \placeholder{x} & \placeholder{x} & \placeholder{x\%} & \placeholder{} \\
\quad Joint vs per-task XGBoost & \placeholder{N} & \placeholder{x} & \placeholder{x} & \placeholder{x} & \placeholder{x\%} & \placeholder{} \\
\quad Joint vs pooled XGBoost & \placeholder{N} & \placeholder{x} & \placeholder{x} & \placeholder{x} & \placeholder{x\%} & \placeholder{} \\
\quad Joint vs XGBoost with registry variables & \placeholder{N} & \placeholder{x} & \placeholder{x} & \placeholder{x} & \placeholder{x\%} & \placeholder{} \\
\bottomrule}
\end{tabular}}
\end{table}

\FloatBarrier
\operasubsection{What to borrow: diseases, outcomes and the population}
Other diseases supplied most of what was worth borrowing. Once outcomes were
shared, adding other diseases changed AUC by \result{P3_auc}{X}. Once diseases
were shared, adding all other outcomes changed it by \result{P4_auc}{Y},
\placeholder{within the prespecified equivalence margin of 0.01 (claim label)}.
The joint model thus predicted all 69 outcomes as well as the disease-pooled
models did one at a time. Borrowing from other outcomes alone, without other
diseases, changed IPA by \result{P2_ipa}{Z} relative to the specialist.

Borrowing from the national population changed the joint model's AUC by
\result{E1_auc}{X} with direct transfer and by \result{E3_auc}{Y} after
adaptation to the cohort. For \opera{}, the corresponding differences were
\result{E2_auc}{X} and \result{E4_auc}{Y} (\cref{fig:encoder-transfer}A).
With 10\%, 25\% and 50\% of the training patients, the adapted national encoder
differed from the cohort encoder by \result{N10_auc}{A}, \result{N25_auc}{B} and
\result{N50_auc}{C} (\cref{fig:encoder-transfer}B). \placeholder{Where the
national encoder helped, by training events and disease}
(\cref{fig:encoder-transfer}C--D).

\begin{figure}[H]
  \resultfigure{fig_external}{National pretraining and adaptation to the cohort}
    {A: from cohort to national pretraining. B: performance with fewer training
    patients. C--D: the gain by training events and by disease.}
  \caption{\textbf{National pretraining, adaptation to the cohort and
  training-set size.} \textbf{A}, One-year AUC of the joint model (filled) and the routed model
  (outline) initialised from the encoder pretrained on the cohort alone, from
  the national encoder (pretrained on the whole population, the cohort
  included) used directly, and from the national encoder adapted to the cohort
  by continued pretraining on its records (95\% CI,
  averaged over the cells evaluated for all six models). Arrows give the paired
  effect of each step, and the bracket the two together. \textbf{B}, Median
  one-year AUC over the same cells of the joint model from the cohort encoder
  and from the adapted national encoder, trained on 10\%, 25\%, 50\% or all of
  the training patients (nested random subsets with the same patients for every
  outcome). Numbers give the median paired difference at each share. The dashed
  line marks the cohort-encoder model trained on all patients, and the open circle the
  share at which the national model reaches it (interpolated on the log scale).
  \textbf{C}, Paired effect of the adapted national versus the cohort encoder
  for each cell, by the number of training events in the cell (line and band:
  median and interquartile range within bins). \textbf{D}, The same effects by
  disease (diamond and bar: median and interquartile range). Tuning and
  test patients were never subsampled.}
  \label{fig:encoder-transfer}
\end{figure}

\FloatBarrier
\operasubsection{How to borrow: own corrections and borrowed corrections}
Letting each prediction weigh its own corrections against borrowed ones improved
modestly on borrowing uniformly. \opera{} changed AUC by \result{R1_auc}{X} and
IPA by \result{R1_ipa}{Y} relative to the joint model, with the higher AUC in
\result{R1_better}{N of M} cells (\resultpct{R1_better}).

In the separately trained controls (\cref{fig:hierarchy}), own disease and
outcome corrections changed AUC by \result{router_own_AB_auc}{X}
(disease-balanced \result{router_own_AB_auc_disease_balanced}{X}), and borrowed
corrections by a further \result{router_borrow_AB_auc}{Y}
(\result{router_borrow_AB_auc_disease_balanced}{Y}). Own corrections thus helped
mainly in the large diseases, which dominate the summary over cells, and
borrowed corrections in the smaller ones, which the disease-balanced summary
weights equally \placeholder{(confirm on test)}. \placeholder{Borrowed
corrections gained most in tasks with few training events}
(\cref{fig:hierarchy}B--C).

Borrowing was broad rather than targeted. \placeholder{The share of routing on
the patient's own disease and outcome, and the effective number of sources.}
Replacing every other disease's correction by its average changed AUC by
\placeholder{X (95\% CI)}, and \placeholder{no single source was indispensable}
(\cref{fig:borrowing-map,fig:supp-borrowing-detail}).

\begin{figure}[H]
  \resultfigure{fig_hierarchy}{Own corrections and borrowing by training support}
    {A: cells by tuning harm. B: own correction and borrowing against training
    events. C: by support band.}
  \caption{\textbf{Own corrections and borrowed corrections, by training
  support.} Three separately trained versions of the routed model are compared:
  shared correction only, with the patient's own disease and outcome
  corrections added (own), and with the corrections of other tasks also
  available (borrowed). Own: own-correction model minus shared-only. Borrowed:
  full model minus own-correction model. Both: full model minus shared-only.
  These are differences between trained models, not an additive split of one
  model's AUC. Gains are AUC pooled over the six horizons within each
  disease--outcome cell. \textbf{A}, Each cell by the joint model's gain over
  the specialist on the tuning set (left of the dashed line, borrowing had
  hurt) against the routed model's gain, coloured by the cell's training events.
  \textbf{B}, Own and borrowed corrections against the target's training events
  at one year (log scale), with a quadratic in log events and a bootstrap band
  over cells, both descriptive. Points are cells with at least 10 training
  events and a 95\% interval narrower than 0.4 (all cells enter the fits).
  \textbf{C}, Median and interquartile range across cells within each band of
  training events, with the number of cells. The spread is between cells, not a
  confidence interval.}
  \label{fig:hierarchy}
\end{figure}

\FloatBarrier
\operasubsection{Borrowing against practical alternatives}
Borrowing within the cohort did not make the sequence models the strongest
predictors.
\opera{} differed from per-task XGBoost by \result{O1_auc}{X} in AUC and
\result{O1_ipa}{Y} in IPA, with the higher AUC in \result{O1_better}{N of M}
cells (\resultpct{O1_better}, \cref{fig:headline,fig:tabular}). XGBoost also gained from
borrowing: pooled across diseases, it differed from per-task XGBoost by
\result{T7_auc}{X}, and the joint model differed from pooled XGBoost by
\result{T6_auc}{Y}. \placeholder{Where the sequence models came closest to or
exceeded XGBoost, by training events and recorded history.}
\placeholder{Calibration slopes and intercepts by model} (\cref{fig:calibration}).
One \opera{} model covered the tasks of \result{narrow_fits}{N} separately
trained specialists.

Borrowing from the national population \placeholder{narrowed / closed / did not
change} the gap to XGBoost. With the adapted national encoder, \opera{} differed from
per-task XGBoost by \result{T5_router_auc}{X} in AUC and \result{T5_router_ipa}{Y}
in IPA, with the higher AUC in \result{T5_router_better}{N of M} cells
(\resultpct{T5_router_better}). The joint model differed by
\result{T5_shared_auc}{Z}. With the national encoder used directly, the
differences were \result{T5E_router_auc}{X} for \opera{} and
\result{T5E_shared_auc}{Z} for the joint model. \placeholder{Where the national encoder narrowed the gap
most, by training events and disease} (\cref{fig:encoder-transfer}).

\begin{figure}[H]
  \resultfigure{fig_headline}{OPERA versus reference methods}
    {Gains over each reference per cell, against XGBoost per cell, against the
    clinical indices, across horizons and by recorded history.}
  \caption{\textbf{Performance of \opera{} relative to reference
  methods.} \textbf{A}, One-year AUC for every available sequence-model and
  tabular variant on the same complete-case task roster. Small points show
  individual task AUCs and diamonds show descriptive means. These are not
  population-pooled AUCs or inferential comparisons. Paired differences with
  confidence intervals are reported in Table 2. \textbf{B}, One-year AUC of \opera{}
  against per-task XGBoost in each cell, with point size reflecting training events.
  \textbf{C}, Share of patients with the event placed in the high-risk group by
  each clinical index and by \opera{} at frozen training-derived cutoffs.
  Specificity was targeted in training and may differ in evaluation. \textbf{D}, Mean AUC gain over per-task XGBoost of
  the joint model, primary OPERA, original-router sensitivity and available
  national OPERA variants at each horizon. \textbf{E}, Improvement in Brier
  score, pooled over horizons, over the specialist and over XGBoost by the number
  of events recorded before treatment ($\times$100, 95\% CI). Within each contrast, models are evaluated on the same test patients.}
  \label{fig:headline}
\end{figure}

\begin{figure}[H]
  \resultfigure{fig_benchmarks}{Comparison with tabular models}
    {A: AUC gains over XGBoost by clinical domain. B: calibration slope.
    C: gain against training events. D--E: reliability.}
  \caption{\textbf{Comparison with tabular models: discrimination and
  calibration.} \textbf{A}, AUC gain over per-task XGBoost, pooled over the six
  horizons, of \opera{}, the joint model and the specialist, and of the joint
  model over XGBoost pooled across diseases and XGBoost given quality-registry
  variables, by clinical domain (median and interquartile range across cells).
  \textbf{B}, Calibration slope of \opera{}, the joint, specialist and XGBoost
  models by domain, over every cell and horizon (ideal 1). \textbf{C}, Gain of \opera{} over XGBoost
  against the target's training events, with the support curve. Points are
  cells with at least 10 training events and a 95\% interval narrower than 0.4
  (all cells enter the curve). \previewonly{Tuning preview: the support curve is a kernel smoother over cells with a 95\% bootstrap band. The hierarchical model is fitted for the test-set analysis.} \textbf{D--E},
  Reliability of \opera{}, the joint model and XGBoost at one year for the
  two example cells, chosen by a rule on training counts (the cell with the most
  training events in its domain among cells with at least 200 training events
  and an event share between 10\% and 50\%). Observed risks are censoring-weighted and point size reflects patients per
  bin. Bins with fewer than 20 patients are not drawn. No model was recalibrated on test data. The full calibration analysis is in \cref{fig:calibration}.}
  \label{fig:tabular}
\end{figure}

\FloatBarrier
\operasubsection{Comparison with the prognostic indices in clinical use}
A single model was compared with each of the \placeholder{N} prognostic indices
in clinical use, on the patients for whom that index was recorded. For
\placeholder{N of N} indices, its size-matched risk groups stratified at least as
well as the index's own categories. The model's lowest group had survival as high
as or higher than the index's lowest category, and its highest group survival as
low as or lower than the index's highest category. \placeholder{For N of these,
the difference in the lowest or highest group was beyond chance. The exceptions
and the comparisons that could not be estimated} (\cref{fig:score-overview}).

In diffuse large B-cell lymphoma, the largest comparison, NCCN-IPI was the main
comparator. Among \placeholder{N} patients with NCCN-IPI recorded, the model's
lowest group, holding as many patients as the NCCN-IPI low-risk category, had
one-year survival of \placeholder{X\% (95\% CI)}, against \placeholder{Y\%} for
that category. In the highest groups, survival was \placeholder{A\% and B\%}
(\cref{fig:risk-groups}).

The same model also covered patients the indices do not. Ten of the 24 diseases
had no dedicated, externally validated index (\cref{tab:index-availability}), and
\placeholder{X\%} of patients had none recorded, whereas the model predicted
every outcome for every patient.

\begin{figure}[H]
  \resultfigure{fig_risk_groups}{Risk groups versus NCCN-IPI}
    {A--B: overall survival by model risk group and by NCCN-IPI group. C:
    one-year overall survival per group. D: cross-classification.}
  \caption{\textbf{Risk groups defined by \opera{} and by NCCN-IPI in
  DLBCL.} Patients with diffuse large B-cell lymphoma and NCCN-IPI recorded at
  diagnosis, compared on overall survival, the endpoint NCCN-IPI was developed
  for. The model places patients into as many groups as NCCN-IPI has
  categories, each with the same number of patients as the corresponding
  NCCN-IPI category, by predicted risk. No outcome information sets the
  boundaries. \textbf{A}, Kaplan--Meier overall survival by model group, with
  numbers at risk below the curves. \textbf{B}, The same by NCCN-IPI category.
  Shaded bands are 95\% bootstrap intervals and the dashed line marks one year.
  \textbf{C}, One-year overall survival per group with 95\% CI, labelled with
  deaths by one year and patients. \textbf{D}, Cross-classification: rows are
  NCCN-IPI categories and columns model groups, each cell with its patients'
  one-year survival and interval above and their number below. Outlined cells
  hold patients placed in the same group by both, cells to the left are patients
  the model places at lower risk, and cells to the right at higher risk.}
  \label{fig:risk-groups}
\end{figure}

\FloatBarrier
\operasubsection{What the predictions rest on, and where they hold}
Hiding the last 7, 14, 21 and 28 days before treatment changed AUC by
\placeholder{A, B, C and D (95\% CIs)}, most in \placeholder{domains}
(\cref{fig:recent-record-gaps}), and only the last year of history retained
\placeholder{X} of the AUC. Laboratory results carried
\placeholder{X} of the attribution to recorded content
(\cref{fig:attribution-main}). \placeholder{Performance by calendar year, region
and patient characteristics, and the prespecified sensitivity analyses}
(\cref{fig:who-benefits,fig:supp-regions}).

\FloatBarrier
\operasection{Discussion}
Borrowing decided whether many prognostic tasks in lymphoid cancers could be
learned at all. Trained on one disease and one outcome, the sequence models
discriminated little better than chance in \placeholder{about half} of the
evaluable tasks. Borrowing from other diseases made most of them useful
predictors (\result{P1_auc}{X}), adding every other outcome cost nothing
measurable, and \opera{} added a smaller gain (\result{R1_auc}{X}). XGBoost
gained far less from pooling (\result{T7_auc}{X}), as it already learned
usefully from few events, and remained the more accurate model
\placeholder{(final \opera{}--XGBoost difference)}. That both model families
gained suggests the shared information is real rather than an artefact of one
architecture.

Classifications separate the lymphoid cancers by biology, and prognostic indices
follow them \citep{alaggio2022who}. Our results suggest that biological
distinctness and predictive separateness are different things. An infection or a
falling neutrophil count may arise through different mechanisms in different
diseases, yet the records that precede it, and how it is investigated and
treated, are similar enough that examples from one disease inform prediction in
another. The gain was largest where a disease had least data of its own, so the
limit for small diseases appears to be the number of examples rather than their
biology. This is not an argument for redrawing diagnostic boundaries. It is an
argument that the boundaries needed for diagnosis and treatment need not also
bound what a predictive model learns from.

Borrowing could also have harmed. Joint training has frequently degraded
individual clinical tasks \citep{mcdermott2020mtl}, and with 69 outcomes of very
different frequency, the rare ones risked being drowned by the common. Instead,
adding every other outcome to a model shared across diseases left discrimination
unchanged within the prespecified margin (\result{P4_auc}{X}). The outcomes arise during the same course of treatment, and the loss balanced
endpoint families, which may both have protected against negative transfer. Other outcomes also added
little on their own: models pooling outcomes within one disease improved far
less than models pooling diseases. What the models borrowed was mainly patients
with the same outcome, not related outcomes in the same patients, and what other
outcomes could add may already have been learned through other diseases. An
equivalent average can hide losses for individual tasks
\placeholder{(number of cells worse, and by how much)}. Still, it means that one
model can cover a whole specialty without the price joint training has often
exacted.

\opera{} was designed to borrow where a task's own data are scarce and to
deviate from the shared pattern where a task differs. It did both, but modestly.
Own corrections helped mainly in the large diseases and borrowed corrections in
the small ones, as partial pooling would predict \citep{EfronMorris1975}, but the
gain over borrowing uniformly was small \placeholder{(revisit with the national
encoder)}. Borrowing was also broad rather than along clinical lines. No single
source was indispensable, and pairs that clinical knowledge suggests should
matter, such as small lymphocytic lymphoma and chronic lymphocytic leukaemia,
showed no specific dependence \placeholder{(inconclusive on tuning, confirm on
test)}. \placeholder{Integrated gradients: features behind each outcome
clinically plausible, diagnosis-specific patterns contributing comparatively
little (confirm).} What transfers appears to be general, such as treatment
exposure, frailty, prior illness and contact with care, and the shared encoder
may already place related diseases together, leaving routing little to add.
Pooling therefore need not be specified by hand. Corrections for disease groups,
such as the indolent B-cell or T-cell lymphomas, would test this directly. These
dependencies describe the fitted model, not biology.

The national arm tests whether foundation models deliver within a specialty.
Population-scale models of health records have been proposed as general
foundations \citep{savcisens2024life2vec,shmatko2025delphi}, but whether they
improve on strong conventional models for a specific clinical question has
rarely been tested \citep{wornow2023shaky}. Here, a model pretrained on the
nationwide Danish registers was fine-tuned for one specialty and compared, on
the same patients, with tuned XGBoost and with models pretrained on the cohort
alone. \placeholder{Interpret the transfer and label-efficiency results, whether
positive, null or harmful.} If national pretraining closes the gap to XGBoost or
reduces the labelled patients needed, it sets a precedent for how national
foundation models can raise performance within specialties: in Denmark, by
pairing the national registers with clinical databases such as DALY-CARE, and
elsewhere wherever national registers exist. If it does not, that is equally
informative, because a well-tuned tabular model on the cohort's own records
would then remain the better choice.

Few settings would allow this study. DALY-CARE holds every Danish adult with a
lymphoid cancer, recorded from the same sources in the same formats, and the
national registers cover everyone in Denmark, so each kind of borrowing could be
measured on the same patients with otherwise identical models. The test set was
opened once, after every choice had been fixed, against a tuned tabular model
and the indices clinicians use. A single model stratified patients at least as well as \placeholder{each of}
the clinical indices and covered the ten diseases without one
\placeholder{(interpret NCCN-IPI and recent-record results)}. These risk groups
describe prognosis under the treatment received, not whether treatment could
safely be changed \citep{vanGeloven2020}, and where XGBoost predicts better, that
accuracy must be weighed against the convenience of one model for every task.
The main limitations concern validation, outcome measurement and precision. Evaluation
was temporal within one health system, and external validation remains
necessary. Test patients' records before 2022 could contribute to
self-supervised pretraining, without outcome labels. Register-derived outcomes
are imperfect, and laboratory outcomes depend on monitoring. Clinical
comparisons are limited to patients with a recorded index. Reported intervals
condition on the fitted models and exclude uncertainty from model selection, and
many uncommon tasks had too few events to be evaluated.

Prognosis in lymphoid cancers need not be learned one disease and one outcome at
a time. For many tasks, borrowing across diseases is what makes learning
possible, and one model can serve a whole specialty without losing accuracy to
negative transfer. \placeholder{If the national encoder helps: specialties need
not build prognosis from their own patients alone. They can start from a model
of the whole population and borrow across their diseases and outcomes, which
sets a precedent for prediction within specialties in Denmark and elsewhere.
If not: the conclusion rests on borrowing within the specialty.} The next steps
are external validation and testing whether acting on the predictions improves
care.

\section*{Data availability}
According to Danish legislation, the health data used in this study may not be
shared publicly. Within the approved DALY-CARE protocol, the data may be
accessed on a collaborative basis under a data processing agreement, and all
analyses must be performed within the DALY-CARE data resource at the Danish
National Genome Center \citep{brieghel2025dalycare}.
\section*{Code availability}
All code, model configurations, outcome definitions and analysis recipes are
publicly available at \placeholder{repository URL and archival DOI}. The
conversion of the registers to the event format is available at
\url{https://github.com/FGA-DIKU/ehr2meds}. Trained model weights were fitted on
patient-level data and are not released.
\section*{Patient and public involvement}
\placeholder{Statement.}
\section*{Acknowledgements}
\placeholder{Funding, collaborators and compute resources.}
\section*{Author contributions}
\placeholder{CRediT statement.}
\section*{Declaration of interests}
\placeholder{Statement.}
